\section{Finite Fields}

\exercise{33.1}{A power of two}
The factorization is $4096=2^{12}$. For every prime $p$ and positive integer $n$, there exists a finite field of order $p^n$. Thus a field with $4096$ elements exists, namely $\operatorname{GF}(2^{12})$.

\exercise{33.2}{An integer with two distinct prime factors}
We have $3127=53\cdot59$. Both factors are prime and distinct, so this is not a prime power. Every finite field has prime-power order; therefore no field of $3127$ elements exists.

\exercise{33.3}{A cube of a prime}
Compute $41^2=1681$ and $1681\cdot41=68{,}921$. Since $41$ is prime, $68{,}921=41^3$ is a prime power, so $\operatorname{GF}(41^3)$ exists and has the required number of elements.

\exercise{33.4}{Primitive eighth roots}
The nonzero elements of $\operatorname{GF}(9)$ form a cyclic group of order $8$. A primitive eighth root is exactly a generator of this group. A cyclic group of order eight has $\varphi(8)=4$ generators, corresponding to powers with exponents $1,3,5,7$ relatively prime to eight. Thus the answer is four.

\exercise{33.5}{Primitive eighteenth roots}
The multiplicative group of $\operatorname{GF}(19)$ is cyclic of order $18$. Its elements of order $18$ are its $\varphi(18)=18(1-1/2)(1-1/3)=6$ generators. Hence there are six primitive eighteenth roots of unity.

\exercise{33.6}{Primitive fifteenth roots}
The multiplicative group of $\operatorname{GF}(31)$ is cyclic of order $30$. It has a unique subgroup of order $15$, whose generators are exactly the elements of order $15$ in the whole group. Their number is $\varphi(15)=15(1-1/3)(1-1/5)=8$. Thus the answer is eight.

\exercise{33.7}{An order absent from the multiplicative group}
The nonzero elements of $\operatorname{GF}(23)$ form a group of order $22$. An element's order must divide $22$, but $10$ does not. Therefore there are no elements of order $10$, and hence zero primitive tenth roots of unity. Merely computing $\varphi(10)$ without checking divisibility would give an incorrect answer.

\exercise{33.8}{True and false statements}
\begin{enumerate}[label=(\alph*)]
\item True. The multiplicative group of a finite field is cyclic.
\item False. The additive group of $\operatorname{GF}(p^n)$ is $(\Z_p)^n$; for $n>1$ every nonzero element has order $p$, so it is not cyclic.
\item True. The complex zeros of $x^{28}-1$ are $e^{2\pi i k/28}$ and form the cyclic group generated by $e^{2\pi i/28}$.
\item False. $60=2^2\cdot3\cdot5$ is not a prime power.
\item True. $125=5^3$ is a prime power.
\item False. $36=2^2\cdot3^2$ is not a prime power.
\item True. The powers $i,i^2=-1,i^3=-i,i^4=1$ show that $i$ has order four.
\item True. Over every finite field there is an irreducible polynomial of each positive degree, including degree $58$ over $\Z_2$.
\item False. Any proposed rational generator has only finitely many primes in its numerator and denominator; its integer powers cannot produce a prime absent from that list. Thus $\Q^*$ is not cyclic.
\item True. An isomorphic image has the same finite size as $F$, and within the fixed algebraic closure there is exactly one subfield of that size, the roots of $x^{|F|}-x$. The image is therefore $F$ itself.
\end{enumerate}

\exercise{33.9}{Two cubic generators inside one algebraic closure}
Both binary cubics $x^3+x^2+1$ and $x^3+x+1$ have value $1$ at $0,1$, so both are irreducible. Thus $\Z_2(\alpha)$ and $\Z_2(\beta)$ each have eight elements. Every element of either field satisfies $x^8-x=0$. In the common algebraic closure this polynomial has at most eight roots, and each of these fields already consists of eight roots. They must therefore be the same set, proving equality of fields rather than just an abstract isomorphism.

\exercise{33.10}{Irreducibles divide a Frobenius polynomial}
Let $f\in\Z_p[x]$ be irreducible of degree $n$, replacing it by a monic associate if necessary. The quotient $E=\Z_p[x]/(f)$ is a field of $p^n$ elements. For $\alpha=x+(f)$, Lagrange's theorem in $E^*$ (and the zero case) gives $\alpha^{p^n}=\alpha$. The minimal polynomial of $\alpha$ is the monic associate of $f$, so it divides every polynomial vanishing at $\alpha$, in particular $x^{p^n}-x$. Hence $f$ divides that polynomial up to its harmless unit coefficient.

\exercise{33.11}{A multiplicative generator has full degree}
Every nonzero element of $F$ is a power of $\alpha$, so belongs to the subfield $\Z_p(\alpha)$. Zero also belongs, hence $\Z_p(\alpha)=F$. If $d=\deg(\alpha,\Z_p)$, the simple extension has $p^d$ elements. Comparing with $|F|=p^n$ gives $d=n$.

\exercise{33.12}{Subfields of a finite field}
Let $E=\operatorname{GF}(p^n)$ and $m\mid n$. Set $K=\{a\in E:a^{p^m}=a\}$. Iterated Frobenius gives $(a+b)^{p^m}=a^{p^m}+b^{p^m}$ and $(ab)^{p^m}=a^{p^m}b^{p^m}$, so $K$ is closed under addition, multiplication, and negatives; inverses of nonzero fixed elements are fixed too. It contains $0,1$, hence is a subfield.
Because $p^m-1$ divides $p^n-1$, the cyclic group $E^*$ has exactly $p^m-1$ solutions of $a^{p^m-1}=1$. Including zero gives $|K|=p^m$.
Any subfield of size $p^m$ consists of roots of $x^{p^m}-x$, so must equal this $K$ by counting. This proves uniqueness. Conversely any subfield has size $p^m$, and the tower law makes $m$ divide $n$.

\exercise{33.13}{Factoring a Frobenius polynomial}
The derivative of $x^{p^n}-x$ is $-1$, so it has no repeated irreducible factor. An irreducible monic polynomial $f$ of degree $d$ divides it exactly when its roots lie in $\operatorname{GF}(p^n)$.
If such a root $a$ lies there, $\Z_p(a)$ is a subfield of order $p^d$, so $d\mid n$ by the tower law. Conversely if $d\mid n$, the field $\Z_p(a)$ of $p^d$ elements embeds as the unique subfield of $\operatorname{GF}(p^n)$ of that size; all its elements satisfy $x^{p^n}=x$, so $f$ divides the polynomial. Thus its monic irreducible factors are exactly those whose degrees divide $n$, each appearing once. Since the polynomial is monic, the product has no additional nonunit constant.

\exercise{33.14}{Euler criterion and a quadratic over seventeen elements}
\textbf{(a)} Choose a generator $g$ of $\Z_p^*$, which has even order $p-1$. Write the nonzero residue $a=g^k$. It is a square exactly when $k$ is even: the squares are the even powers of $g$. Also
\[a^{(p-1)/2}=1\quad\Longleftrightarrow\quad(p-1)\mid k(p-1)/2\quad\Longleftrightarrow\quad2\mid k.\]
This proves the criterion, including both implications.
\textbf{(b)} Modulo seventeen, $6^2\equiv2$, $6^4\equiv4$, and $6^8\equiv16=-1\ne1$. Hence six is not a quadratic residue, so $x^2-6$ has no root in $\Z_{17}$ and is irreducible there.
